% =====================================================================
\section{Luz y semiconductores}
% =====================================================================

Los dispositivos optoelectrónicos (fotoconductores, fotodiodos, celdas solares, LEDs y
láseres) son dispositivos semiconductores diseñados para \emph{absorber} o \emph{emitir}
luz de forma eficiente. Todos se basan en la misma idea: un fotón con energía suficiente
puede llevar un electrón de la banda de valencia a la de conducción (\emph{absorción},
genera un par electrón--hueco), y un electrón que cae de la banda de conducción a la de
valencia puede liberar esa energía como un fotón (\emph{emisión}).

\subsection{El fotón: energía y longitud de onda}

La luz está formada por fotones de energía
\begin{equation}
  \Eph = h\nu = \frac{hc}{\lambda},
\end{equation}
con $h=\SI{6.626e-34}{J.s}=\SI{4.136e-15}{\eV.s}$ y $c=\SI{3e8}{m/s}$. Como
$hc = \SI{1.24}{\eV.\micro m}$, en unidades prácticas:

\begin{clave}[Relación energía -- longitud de onda \enhoja]
\[
  \lambda\,[\si{\micro m}] = \frac{1.24}{\Eph\,[\si{\eV}]}
  \qquad\Longrightarrow\qquad
  \lG\,[\si{\micro m}] = \frac{1.24}{\EG\,[\si{\eV}]}
\]
$\lG$ es la \textbf{longitud de onda de corte} del semiconductor: la longitud de onda
cuyo fotón tiene exactamente la energía del \emph{gap}. Fotones con $\lambda<\lG$
($\Eph>\EG$) pueden ser absorbidos; fotones con $\lambda>\lG$ ($\Eph<\EG$) atraviesan el
material: el semiconductor es \textbf{transparente} para ellos.
\end{clave}

\begin{figure}[H]
\centering
\begin{tikzpicture}[x=7.2cm,y=1cm]
  % banda visible con colores reales (modelo "wave" de xcolor, en nm)
  \foreach \l in {380,384,...,748}{
    \definecolor{tmpc}{wave}{\l}
    \fill[tmpc] ({(\l/1000-0.2)},0) rectangle ({(\l/1000-0.2)+0.0042},0.45);
  }
  \fill[violet!25] (0,0) rectangle (0.18,0.45);
  \fill[red!15] (0.552,0) rectangle (1.8,0.45);
  \node[font=\small] at (0.09,0.225) {UV};
  \node[font=\small] at (1.2,0.225) {infrarrojo (IR)};
  \draw[flecha] (0,0) -- (1.85,0) node[right] {$\lambda$ (\si{\micro m})};
  \foreach \x in {0.2,0.4,...,2.0}{
    \draw ({\x-0.2},0) -- ({\x-0.2},-0.1) node[below,font=\scriptsize] {\pgfmathprintnumber[fixed,precision=1]{\x}};
  }
  \node[font=\scriptsize] at (0.37,-0.55) {visible};
  % materiales: lambda_G
  \foreach \lam/\nom/\h in {0.365/GaN/1.0, 0.413/SiC/1.6, 0.46/ZnSe/2.2, 0.549/GaP/1.0,
                             0.873/GaAs/1.0, 0.919/InP/1.6, 1.107/Si/1.0, 1.65/InGaAs/1.0, 1.88/Ge/1.6}{
    \draw[gris,thin] ({\lam-0.2},0.45) -- ({\lam-0.2},\h);
    \node[above,font=\scriptsize,inner sep=1pt] at ({\lam-0.2},\h) {\nom};
  }
  % ventanas de fibra
  \foreach \lam in {1.31,1.55}{
    \draw[azul,densely dashed] ({\lam-0.2},0) -- ({\lam-0.2},-0.9);
  }
  \node[azul,font=\scriptsize,align=center] at (1.23,-1.15) {ventanas de fibra óptica\\1.31 y \SI{1.55}{\micro m}};
  \node[font=\scriptsize,align=left,anchor=west] at (0.95,2.1) {marcas: $\lG=1.24/\EG$ de cada material};
\end{tikzpicture}
\caption{Espectro óptico y longitud de onda de corte $\lG$ de semiconductores comunes (300\,K).
El rango visible va de \SIrange{0.38}{0.75}{\micro m}, es decir de \SIrange{3.3}{1.65}{\eV}.}
\label{fig:espectro}
\end{figure}

\begin{table}[H]
\centering\small
\caption{Gap de algunos semiconductores a 300\,K. La columna ``tipo'' es fundamental: decide si
el material sirve para emitir luz (sección~\ref{sec:led}).}
\label{tab:gaps}
\begin{tabular}{lcccl}
\toprule
Material & $\EG$ (\si{\eV}) & $\lG$ (\si{\micro m}) & Tipo de gap & Uso típico \\
\midrule
Ge                       & 0.66 & 1.88 & indirecto & detectores IR \\
In$_{0.53}$Ga$_{0.47}$As & 0.75 & 1.65 & directo   & detectores para fibra óptica \\
Si                       & 1.12 & 1.11 & indirecto & fotodiodos, celdas solares \\
InP                      & 1.35 & 0.92 & directo   & sustrato, ventanas, láseres \\
GaAs                     & 1.42 & 0.87 & directo   & LEDs IR, láseres, celdas III-V \\
GaP                      & 2.26 & 0.55 & indirecto & LEDs (dopado con N) \\
SiC (6H)                 & 3.0  & 0.41 & indirecto & LEDs azules (antiguos) \\
GaN                      & 3.4  & 0.36 & directo   & LEDs azules/blancos (con InGaN) \\
\bottomrule
\end{tabular}
\end{table}

\begin{ojo}[gap del GaP \diapo{6}]
La tabla de la diapositiva pinta al GaP (\SI{2.3}{\eV}) de verde como si fuera un emisor
natural de verde. El GaP es de \textbf{gap indirecto}: sólo emite con eficiencia razonable
gracias a impurezas \emph{isoelectrónicas} (N, o el par Zn--O) que se agregan a propósito
(ver sección~\ref{sec:led}). Los valores de la diapo están redondeados:
Ge \SI{0.66}{\eV}, Si \SI{1.12}{\eV}, GaAs \SI{1.42}{\eV}.
\end{ojo}

\subsection{Absorción en semiconductores de gap directo e indirecto}

En la absorción hay que conservar la \textbf{energía} y el \textbf{momento cristalino}
($\hbar k$). El momento de un fotón es despreciable: $k_{\text{fotón}}=2\pi/\lambda\sim
\SI{e5}{cm^{-1}}$, mientras que la zona de Brillouin se extiende hasta $k\sim\pi/a\sim
\SI{e8}{cm^{-1}}$. Por eso, en un diagrama $E$--$k$, \emph{la transición debida sólo al
fotón es vertical}.

\begin{figure}[H]
\centering
\begin{tikzpicture}[scale=0.95]
  % ---------- (a) gap directo ----------
  \begin{scope}
    \draw[flecha] (-2,0) -- (2.2,0) node[right] {$k$};
    \draw[flecha] (0,-1.8) -- (0,3.2) node[above] {$E$};
    \draw[bandac,domain=-1.5:1.5,smooth] plot (\x,{1.4+0.6*\x*\x});
    \draw[bandav,domain=-1.6:1.6,smooth] plot (\x,{-0.15-0.45*\x*\x});
    \node[colEc,right] at (1.5,2.75) {BC};
    \node[colEv,right] at (1.6,-1.3) {BV};
    \node[elec] (e) at (0,1.4) {};
    \node[hueco] (h) at (0,-0.15) {};
    \draw[flecha,colFoton,line width=1pt] (0,-0.05) -- (0,1.28);
    \draw[foton] (-1.7,0.9) -- (-0.15,0.65);
    \node[colFoton,font=\small] at (-1.3,1.3) {$h\nu$};
    \draw[<->,gris] (0.35,-0.15) -- node[right,font=\small] {$\EG$} (0.35,1.4);
    \node[font=\small,align=center] at (0,-2.4) {(a) gap \textbf{directo}\\(GaAs, InP, GaN)};
  \end{scope}
  % ---------- (b) gap indirecto ----------
  \begin{scope}[xshift=6.3cm]
    \draw[flecha] (-2,0) -- (2.4,0) node[right] {$k$};
    \draw[flecha] (0,-1.8) -- (0,3.2) node[above] {$E$};
    \draw[bandac,domain=-1.6:2.3,smooth,samples=60] plot (\x,{1.1+0.75*(\x-1.2)*(\x-1.2)/(1+0.35*(\x-1.2)*(\x-1.2))});
    \draw[bandav,domain=-1.6:1.6,smooth] plot (\x,{-0.15-0.45*\x*\x});
    \node[colEc,right] at (2.3,1.6) {BC};
    \node[colEv,right] at (1.6,-1.3) {BV};
    \node[hueco] at (0,-0.15) {};
    \node[elec] at (1.2,1.1) {};
    % foton vertical hasta estado virtual
    \draw[flecha,colFoton,line width=1pt] (0,-0.05) -- (0,0.95);
    \draw[gris,dotted] (0,1.05) -- (1.1,1.05);
    \draw[flecha,violeta,line width=0.9pt] (0.05,1.05) -- (1.08,1.1);
    \node[violeta,font=\small] at (0.75,0.65) {fonón};
    \draw[<->,gris] (1.75,-0.15) -- node[right,font=\small] {$\EG$} (1.75,1.1);
    \draw[gris,dotted] (0,-0.15) -- (1.8,-0.15);
    \node[font=\small,align=center] at (0.2,-2.4) {(b) gap \textbf{indirecto}\\(Si, Ge, GaP)};
  \end{scope}
\end{tikzpicture}
\caption{Absorción de un fotón en los diagramas $E$--$k$. En (a) el mínimo de la BC está
justo encima del máximo de la BV y basta el fotón. En (b) el electrón tiene que cambiar de
$k$: necesita además \emph{absorber o emitir un fonón} (vibración de la red), que aporta
el momento faltante (y una pequeña energía $E_{ph}\sim\SI{50}{meV}$).}
\label{fig:Ek}
\end{figure}

\begin{intuicion}[por qué importa el tipo de gap]
Un proceso que necesita fotón \emph{y} fonón a la vez es mucho menos probable que uno que
necesita sólo el fotón. Consecuencias:
\begin{itemize}
  \item \textbf{Absorción:} en un material indirecto el coeficiente de absorción $\alpha$
  crece lentamente al pasar $\lG$; hacen falta espesores grandes (en Si, una celda solar
  necesita $\sim\SI{100}{\micro m}$ o más; en GaAs alcanzan unos pocos \si{\micro m}).
  \item \textbf{Emisión:} es el mismo argumento al revés. En Si la recombinación
  banda a banda radiativa es tan lenta que gana la recombinación por trampas, que
  libera la energía como calor. Por eso \emph{no hay LEDs de silicio}.
\end{itemize}
Con absorción de fonón se pueden absorber fotones con $h\nu=\EG-E_{ph}$; con emisión de
fonón, $h\nu=\EG+E_{ph}$ (las tres variantes del diagrama de la \diapo{3}).
\end{intuicion}

\subsection{Coeficiente de absorción y ley de Beer--Lambert}

Supongamos un haz de luz monocromática que entra a un semiconductor por $x=0$. Cada
unidad de longitud recorrida absorbe una fracción fija $\alpha$ de los fotones
que quedan:
\[
  \frac{\dd\Phi}{\dd x} = -\alpha\,\Phi
  \qquad\Longrightarrow\qquad
  \Phi(x) = \Phi_0\,e^{-\alpha x},\qquad I(x)=I_0\,e^{-\alpha x},
\]
donde $\Phi$ es el flujo de fotones (fotones/\si{cm^2.s}) e $I$ la intensidad
(\si{W/cm^2}). $\alpha$ [\si{cm^{-1}}] es el \textbf{coeficiente de absorción}, y
$1/\alpha$ es la \textbf{profundidad de penetración}: la distancia en que la luz cae a
$1/e\approx\SI{37}{\percent}$. A una profundidad de $3/\alpha$ ya se absorbió el
\SI{95}{\percent}.

\begin{figure}[H]
\centering
\begin{tikzpicture}
\begin{semilogyaxis}[
  width=0.62\textwidth, height=6.4cm,
  xmin=0.3, xmax=1.9, ymin=1, ymax=2e6,
  xlabel={$\lambda$ (\si{\micro m})}, ylabel={$\alpha$ (\si{cm^{-1}})},
  grid=both, minor grid style={gray!15}, major grid style={gray!30},
  legend pos=south west, legend style={font=\small},
  every axis plot/.append style={line width=1.1pt,smooth},
  axis y line*=left,
]
\addplot[azul] coordinates {(0.30,1.7e6)(0.35,1.0e6)(0.40,9.5e4)(0.45,2.6e4)(0.50,1.1e4)(0.55,6.4e3)(0.60,4.1e3)
  (0.65,2.8e3)(0.70,1.9e3)(0.75,1.3e3)(0.80,850)(0.85,535)(0.90,306)(0.95,157)(1.00,64)(1.05,16)(1.10,3.5)(1.12,1.2)};
\addlegendentry{Si}
\addplot[rojo] coordinates {(0.30,1.5e6)(0.40,7e5)(0.50,1.2e5)(0.60,5e4)(0.70,2.8e4)(0.80,1.3e4)(0.85,8e3)
  (0.865,3e3)(0.875,4e2)(0.885,30)(0.89,1)};
\addlegendentry{GaAs}
\addplot[verde] coordinates {(0.40,5e5)(0.50,2e5)(0.60,8e4)(0.70,5e4)(0.80,3e4)(0.90,2e4)(0.915,6e3)(0.925,5e2)(0.935,10)(0.94,1)};
\addlegendentry{InP}
\addplot[naranja] coordinates {(0.40,6e5)(0.60,2e5)(0.80,6e4)(1.00,2e4)(1.20,1e4)(1.40,7e3)(1.50,4e3)(1.55,7e2)
  (1.60,1.5e2)(1.70,30)(1.80,5)(1.86,1)};
\addlegendentry{Ge}
\end{semilogyaxis}
\begin{semilogyaxis}[
  width=0.62\textwidth, height=6.4cm,
  xmin=0.3, xmax=1.9, ymin=1, ymax=2e6,
  axis x line=none, axis y line*=right,
  ytick={1e-4,1e-3,1e-2,1e-1,1,10,100,1000,10000,100000,1000000},
  yticklabels={},
  extra y ticks={1e2,1e3,1e4,1e5},
  extra y tick labels={\SI{100}{\micro m},\SI{10}{\micro m},\SI{1}{\micro m},\SI{0.1}{\micro m}},
  extra y tick style={tick label style={font=\scriptsize}},
  ylabel={$1/\alpha$}, ylabel style={font=\small},
]
\end{semilogyaxis}
\end{tikzpicture}
\caption{Coeficiente de absorción en función de $\lambda$ (valores aproximados a 300\,K,
adaptado de Pierret, fig.~3.20). A la derecha, la profundidad de penetración $1/\alpha$
correspondiente. Notar el corte \emph{abrupto} de los materiales directos (GaAs, InP) en
$\lG$ y el corte \emph{gradual} del Si (indirecto). El Ge tiene un gap indirecto de
\SI{0.66}{\eV} y uno directo de \SI{0.80}{\eV} ($\approx\SI{1.55}{\micro m}$): de ahí el
escalón.}
\label{fig:alfa}
\end{figure}

\begin{table}[H]
\centering\small
\caption{Profundidad de penetración $1/\alpha$ en silicio. La luz azul se absorbe en la
superficie; la infrarroja cercana penetra cientos de micrones.}
\label{tab:penetracion}
\begin{tabular}{lccccccc}
\toprule
$\lambda$ (\si{\micro m}) & 0.40 & 0.50 & 0.60 & 0.70 & 0.80 & 0.90 & 1.00 \\
\midrule
$\alpha$ (\si{cm^{-1}}) & \num{9.5e4} & \num{1.1e4} & \num{4.1e3} & \num{1.9e3} & 850 & 306 & 64 \\
$1/\alpha$ (\si{\micro m}) & 0.1 & 0.9 & 2.4 & 5.3 & 12 & 33 & 156 \\
\bottomrule
\end{tabular}
\end{table}

\subsection{Fotogeneración}
\label{sec:fotogeneracion}

Cada fotón absorbido (con $\Eph>\EG$) genera un par electrón--hueco. La tasa de
fotogeneración por unidad de volumen es entonces el número de fotones absorbidos por
unidad de volumen y de tiempo, $-\dd\Phi/\dd x$:

\begin{clave}[Tasa de fotogeneración \enhoja]
\[
  G_L(x,\lambda) = \left.\frac{\partial n}{\partial t}\right|_{\text{luz}}
  = \left.\frac{\partial p}{\partial t}\right|_{\text{luz}}
  = \alpha\,\Phi_0\,e^{-\alpha x} \equiv G_{L0}\,e^{-\alpha x}
  \qquad [\si{cm^{-3}.s^{-1}}]
\]
con $\Phi_0 = \dfrac{(1-R)\,P}{h\nu\,A}$ el flujo de fotones que \emph{entra} al
semiconductor (potencia $P$ sobre área $A$, reflectividad $R$). Electrones y huecos se
generan \textbf{de a pares}: $\Delta n=\Delta p$ en el instante de la generación.
\end{clave}

\begin{figure}[H]
\centering
\begin{tikzpicture}
  % bloque semiconductor
  \fill[azulclaro] (0,0) rectangle (6,1.4);
  \draw (0,0) rectangle (6,1.4);
  \node[font=\small,anchor=south east] at (5.95,0.05) {semiconductor};
  % sombreado de intensidad decreciente
  \foreach \i in {0,...,29}{
    \pgfmathsetmacro{\op}{0.85*exp(-\i/9)}
    \fill[naranja,opacity=\op] ({\i*0.2},0) rectangle ({\i*0.2+0.2},1.4);
  }
  \foreach \y in {0.3,0.7,1.1}{
    \draw[foton] (-1.6,\y) -- (-0.05,\y);
  }
  \node[colFoton,above] at (-0.9,1.4) {$h\nu$, $\lambda$};
  % grafico debajo
  \begin{scope}[yshift=-3.2cm]
    \draw[flecha] (0,0) -- (6.4,0) node[right] {$x$};
    \draw[flecha] (0,0) -- (0,2.4) node[above] {$G_L(x)$};
    \draw[naranja,line width=1.1pt,domain=0:6,smooth,samples=60] plot (\x,{2*exp(-\x/1.8)});
    \draw[gris,dashed] (1.8,0) node[below,black] {$1/\alpha$} -- (1.8,{2*exp(-1)});
    \draw[gris,dashed] (0,{2*exp(-1)}) node[left,black] {$G_{L0}/e$} -- (1.8,{2*exp(-1)});
    \node[left] at (0,2) {$G_{L0}$};
    \node[font=\small,align=left,anchor=west] at (3.0,1.4) {$G_L=G_{L0}\,e^{-\alpha x}$\\$\alpha$ grande: todo en la superficie\\$\alpha$ chico: casi uniforme};
  \end{scope}
\end{tikzpicture}
\caption{La luz se atenúa exponencialmente al penetrar y con ella la tasa de generación de
pares \diapo{4}. Si el dispositivo es mucho más delgado que $1/\alpha$, se puede aproximar
$G_L\approx$ constante, que es la hipótesis que usa Pierret en las deducciones de
fotodiodos.}
\label{fig:fotogeneracion}
\end{figure}

\paragraph{Reflexión en la superficie.} Al pasar del aire ($n_1=1$) a un semiconductor de
índice $n$ con incidencia normal se refleja una fracción
\[
  R = \left(\frac{n-1}{n+1}\right)^2 .
\]
Para Si ($n\approx3.5$) da $R\approx\SI{31}{\percent}$: ¡casi un tercio de la luz se pierde
antes de entrar! Por eso todos los fotodetectores y celdas llevan \textbf{capa
antirreflectante}: una película de índice $n_{AR}\approx\sqrt{n}$ y espesor
$\lambda/(4n_{AR})$ hace que las ondas reflejadas en sus dos caras se cancelen. Las celdas
solares además \emph{texturizan} la superficie (pirámides) para que la luz rebotada
vuelva a incidir sobre el semiconductor.

\begin{ejemplo}[¿qué material para qué luz?]
\textbf{a)} ¿Puede un fotodiodo de Si detectar la luz de un láser de fibra de
\SI{1.55}{\micro m}? \quad
$\Eph = 1.24/1.55 = \SI{0.80}{\eV} < \EG(\text{Si}) = \SI{1.12}{\eV}$. \textbf{No}: el Si
es transparente a esa luz. Sirven Ge ($\lG=\SI{1.88}{\micro m}$) o
In$_{0.53}$Ga$_{0.47}$As ($\lG=\SI{1.65}{\micro m}$).

\textbf{b)} Luz roja de \SI{0.65}{\micro m} incide sobre Si. ¿Qué espesor absorbe el
\SI{90}{\percent} de lo que entra? \quad Interpolando la tabla~\ref{tab:penetracion},
$\alpha\approx\SI{2.8e3}{cm^{-1}}$. Se pide $1-e^{-\alpha d}=0.9\Rightarrow
d=\ln 10/\alpha = 2.3/\SI{2.8e3}{cm^{-1}}\approx\SI{8}{\micro m}$.
\end{ejemplo}

\begin{ojo}[\diapo{5}]
La frase ``el fotodiodo PN tiene una respuesta en frecuencia limitada a aprox.\
\SI{10}{MHz}'' aparece en la diapositiva del coeficiente de absorción pero no tiene que ver
con $\alpha$ directamente: se explica en la sección~\ref{sec:pn-velocidad} (es la lentitud
de la difusión de los portadores generados fuera de la zona de carga espacial).
\end{ojo}
